\chapter{Extending applications and maintaining the code}\label{ch:development}
\section{Begin with a physical and software contract}
State the target equation, units, frame, initial/boundary conditions, source
normalization, valid parameter domain and output observable before installing
a callback. Distinguish a scientific assumption from an implementation detail.
If a required physical closure is unspecified, request that decision rather
than embedding a guessed default.

Choose the owner of each responsibility. Generic algorithms belong in the
core/shared model; target-specific geometry, fit coefficients and observation
operators belong in applications. Keep application adapters small enough to
review without reading all particle algorithms. Preserve the null/default
path when adding an optional hook.

\section{A practical new-application sequence}
\begin{enumerate}
\item Start from a small maintained native test or the nearest application
with the same geometry and storage needs. Do not copy unrelated physics.
\item Add an explicitly named compile-time deck under \path{input}; select
its source directory, species, dimension, layouts and supported modules.
\item Implement a runtime owner using the generic grammar where appropriate,
with unknown-key rejection and cross-field validation.
\item Register node/particle storage and callbacks before layout freezing.
Build/allocate the mesh and initialize complete owner state.
\item Add sources with a physical-rate ledger, the correct event law and
native ownership. Derive time steps and weights from the declared numerical
policy and measure the realized values.
\item Install the mover/boundary/reaction callbacks with exact lifecycle and
list-ownership rules. Complete exchange before dependent sampling/output.
\item Add independent model tests, a bounded native end-to-end case and
negative setup controls before a costly science campaign.
\item Document exact commands, files, output interpretation, requirements
and unsupported paths. Register application tests globally only when the
adapter can truthfully import every owned result.
\end{enumerate}

\section{Comments and documentation}
Explain equations, dimensional factors, sign conventions, approximations,
algorithmic order, tolerances and typed failure behavior. A comment repeating
the next assignment does not explain why the assignment is physically correct.
Document owner/ghost and generation boundaries whenever state crosses MPI.
Keep nearby READMEs current and update this manual when a public workflow or
capability changes.

Public interfaces should validate inputs before mutating state. Prefer an
explicit unsupported/invalid outcome to a plausible fallback value.
Separate model rejection from missing-data and build failures. Never clamp
an invalid physical parameter solely to force a downstream solver to continue.
When numerical roundoff bounds are appropriate, derive and document them
separately from physical acceptance tolerances.

\section{Portability and reproducibility}
Use configured argv-based compilers/launchers; avoid host-specific paths in
test implementations. Keep site settings in the required machine profiles
or established native configuration workflow. Test both packed/aligned
particle layouts when touching record access; verify DIM1/2/3 when touching
tree or face indexing. Exercise one-/multi-rank ownership and changed
decomposition when touching exchange or persistence.

Compiler sanitizers cover only instrumented code. Linking an uninstrumented
archive into a sanitized consumer does not instrument that archive.
An optional CUDA compile probe does not qualify a device physics algorithm.
Retain compiler flags, binary hashes, input/data fingerprints and exact
commands with each qualification receipt.

\section{Performance without changing the experiment}
Profile initialization, transport, collisions, chemistry, exchange, sampling
and output separately. Record mesh size, particle population and ranks along
with wall time. CPU time in one rank can be synchronization wait rather than
local physical work. Optimize allocation/search structure only after proving
equivalent rate, random-draw and ownership behavior where those are required.

Do not reduce a refinement matrix, filter a shock branch or increase a frozen
routine timeout and then report the original acceptance as passed.
Keep performance diagnostics distinct from physical qualification.
Write scalable observables and avoid unnecessary full-volume text output.

\section{Source and reference preservation}
Use read-only Git inspection under the current repository policy. Leave edits
unstaged and uncommitted; do not alter repository metadata or Git control
files. A historical document suggesting commits is not authorization to run
them. Import only reviewed source/reference inputs; generated run logs,
plots and output trees are evidence rather than application source.
Before every native rebuild, preserve the root/idle, generated-build removal,
configuration and top-level -j16 sequence from Chapter~\ref{ch:installation}.
\source{src/README.md; src/test/README.md}
